\subsection{Degrees forty, forty-four, and forty-eight}
\label{sec:k40-k44-k48}

\begin{theorem}[K17-K40-K44-K48]\label{thm:k40-k44-k48}
The stipulated arithmetic theorem of Alp\"oge implies
\[
               K_{40}(\mathbf Z)=K_{44}(\mathbf Z)
                               =K_{48}(\mathbf Z)=0.
 \tag{HK3.1}
\]
Together with the preceding results,
\[
                K_{4k}(\mathbf Z)=0\qquad(1\leq k\leq12).
 \tag{HK3.2}
\]
For every odd prime $p$ and every $k$ in this range,
$A_p^{[-2k]}=0$ and $\mathcal S_p(2k)=0$.
The full-domain Soul\'e characters $\kappa_m$ are surjective
for every odd $m$ with $3\leq m\leq25$, at every odd prime $p$.
\end{theorem}

\begin{proof}
Fix $e\in\{20,22,24\}$, so the requested $K$-degree is $2e$.
For an odd prime $p<2^{31}$, the cited Hart--Harvey--Ong
Vandiver verification and (K2)--(K3) give
$K_{2e}(\mathbf Z)\{p\}=0$. The 2-primary part is zero by
the cited Rognes--Weibel theorem. It remains to treat
primes $p>2^{31}$; there is no prime equal to $2^{31}$.

First suppose that $p$ divides none of the residual integers
$M_{e,i}$ in (HH3.3). The actual datum
$\mathrm{GD}(\Omega,\mathbf Q(\sqrt{-3}),p,e,1)$, with
modulus two, is valid as in the preceding general coverage
theorem: $e$ is positive and even, $p\ne2,3$, and $p-1>2e$.
The exact combined gcd (HH3.4) has no prime factor outside
the displayed small factors and those residual integers.
Consequently at least one of $N_{e,7}$ and $N_{e,13}$ is a
$p$-adic unit. Choose that auxiliary prime $r=7$ or $13$.
It splits in $\mathbf Q(\sqrt{-3})$ and is different from $p$.
The complete full-space models of
Proposition~\ref{prop:hecke-higher3}, the uniform rank-three
coverage and rank-two conductor theorem, and
Corollary~\ref{cor:general-hecke-unit} give
$\mathcal S_p(e)=0$. By (K5),
$K_{2e}(\mathbf Z)\{p\}=0$.
The gcd was taken after multiplying the two branch values
at each $r$, so the selected $r$ covers both families
simultaneously.

Now suppose that $p\mid M_{e,i}$ for some displayed residual
integer. Proposition~\ref{prop:u3-minus7-higher3}, through the
full residue-ring criterion of Lemma~\ref{lem:u3-residual-ring},
makes the entire compact comparison over
$\mathbf Q(\sqrt{-7})$ a unit at every coefficient place
above this $p$. Proposition~\ref{prop:minus7-gl2-higher3}
and its full integral basis likewise make the complete
cubic-twist classical comparison a unit at every such place.
These statements apply to every prime divisor of each
$M_{e,i}$, irrespective of its primality or squarefreeness.
Since $p>2^{31}>\max(2e+1,7)$,
Proposition~\ref{prop:minus7-general-source} applies and gives
$\mathcal S_p(e)=0$ and the same primary $K$-group vanishing.
The full classical comparison includes all odd nebentypus
characters and all three untwisting values at $2$; original
conductor-$49$ systems are retained.

Every prime has now been treated. The positive $K$-groups of
$\mathbf Z$ in degrees divisible by four are finite, so
vanishing of all their primary parts proves (HK3.1).
The previously proved range $1\leq k\leq9$ supplies (HK3.2).
The exact class-component equivalence (K2) and strict duality
(K5) give the two cohomological assertions at every odd $p$.
Finally for odd $m=2k+1$ the full-domain index formula (K18)
has $2m-2=4k$ and gives
\[
 [\mathbf Z_p(m):\kappa_m(N_p)]
          =|K_{4k}(\mathbf Z)\{p\}|=1.
\]
This proves the stated full-domain surjectivity in the exact
cyclotomic-product normalization (K19). It does not assert
unit values on a deeper Galois filtration subgroup.
\end{proof}

\begin{corollary}[K17-TWELVE-DEGREES-PERIOD]
\label{cor:twelve-degrees-period}
For every odd prime $p$ and positive integer $k$ whose residue
modulo $(p-1)/2$ has a representative in
\[
                             -16\leq j\leq12,
 \tag{HK3.3}
\]
one has $K_{4k}(\mathbf Z)\{p\}=0$.
In the cyclotomic tower of
Proposition~\ref{prop:componentwise-iwasawa-control}, one has
\[
              A_n^{[-2k]}=X^{[-2k]}=0\qquad(n\geq0).
\]
For $p\geq5$, the actual ramified component $M^{[2k+1]}$
is integrally free of rank one over $\Lambda$.
\end{corollary}

\begin{proof}
For $1\leq j\leq12$, (HK3.2) and (K2) give
$A_0^{[-2j]}=0$; for $j=0$ use the trivial class component.
For $-16\leq j\leq-1$, use
Proposition~\ref{prop:positive-iwasawa-components}.
Congruence modulo $(p-1)/2$ gives the same exponent $-2k$
modulo $p-1$. By (K2), this gives the claimed primary
$K$-group vanishing. The componentwise control proposition therefore
gives all finite-layer and inverse-limit vanishings, and its
integral freeness implication gives $M^{[2k+1]}$.
When that odd character is $\omega$ or $\omega^{-1}$, use
the exceptional Coleman calculation in
Corollary~\ref{cor:odd-companion-exclusion}. These cases
exhaust all residues and prove the assertion.
\end{proof}

The entire range (HK3.2) contains ten new specified vanishing
degrees beyond the known degrees four and eight. The finite
character interval (HK3.3) need not exhaust all characters at
a large prime. Neither all-prime Vandiver nor the remaining
all-prime tripod-kernel assertion is inferred from it.
